% Classe appositamente creata per tesi di Ingegneria Informatica all'università Roma Tre
\documentclass{TesiDiaUniroma3}
\usepackage{algorithm}
\usepackage{algpseudocode}
\usepackage{float} 
\usepackage{eurosym}
\usepackage[algo2e]{algorithm2e} %aggiunto sofia
% \usepackage[ruled,vlined]{algorithm2e} %aggiunto sofia
\usepackage{tikz} %aggiunto sofia per cap 2
\usetikzlibrary{shapes.geometric, arrows, positioning} %aggiunto sofia per cap 2
\usepackage{longtable} %aggiunti sofia
\usepackage{array} %aggiunti sofia
\usepackage{amssymb} %aggiunto sofia
\usepackage{amsmath} %aggiunto sofia
\usepackage{graphicx}

% --- INIZIO dati relativi al template TesiDiaUniroma3
% dati obbligatori, necessari al frontespizio
\titolo{Simulazione dei Processi Industriali e Logistici}
\autore{Sofia Garofoli \\  Davide Tommasino}
\relatore{Prof. Ludovica Adacher}
\correlatore{Prof. Nome Cognome} % modifica anche TesiDiaUniroma3.cls se vuoi avere un correlatore
\annoAccademico{2025/2026}

% dati opzionali
\dedica{ a me} % solo se nel documento si usa il comando \generaDedica
% --- FINE dati relativi al template TesiDiaUniroma3

% --- INIZIO richiamo di pacchetti di utilità. Questi sono un esempio e non sono strettamente necessari al modello per la tesi.
\usepackage[plainpages=false]{hyperref}	% generazione di collegamenti ipertestuali su indice e riferimenti
\usepackage[all]{hypcap} % per far si che i link ipertestuali alle immagini puntino all'inizio delle stesse e non alla didascalia sottostante
\usepackage{amsthm}	% per definizioni e teoremi
\usepackage{amsmath}	% per ``cases'' environment
% --- FINE riachiamo di pacchetti di utilità

\begin{document}
% ----- Pagine di fronespizio, numerate in romano (i,ii,iii,iv...) (obbligatorio)
\frontmatter
\generaFrontespizio
% \generaDedica
% \ringraziamenti{ringraziamenti}	% inserisce i ringraziamenti e li prende in questo caso da ringraziamenti.tex
\introduzione{introduzione}		% inserisce l'introduzione e la prende in questo caso da introduzione.tex
\generaIndice
\generaIndiceFigure

% ----- Pagine di tesi, numerate in arabo (1,2,3,4,...) (obbligatorio)
\mainmatter
% il comando ``capitolo'' ha come parametri:
% 1) il titolo del capitolo
% 2) il nome del file tex (senza estensione) che contiene il capitolo. Tale nome \`e usato anche come label del capitolo
\capitolo{Formulazione del problema}{capitolo-1}
\capitolo{Algoritmi Euristici per il Problema}{capitolo-2}
\capitolo{Analisi sulle Istanze del Problema}{capitolo-3}



% ----- Parte finale della tesi (obbligatorio)
\backmatter
\conclusioni{conclusioni}

% Bibliografia con BibTeX (obbligatoria)
% Non si deve specificare lo stile della bibliografia
% NON CREDO SERVA
% \bibliography{bibliografia} % inserisce la bibliografia e la prende in questo caso da bibliografia.bib

\end{document}
